\chapter{Evaluation and Results}
\label{ch:evaluation_results}

\label{sec:evaluation_introduction}

To evaluate the effectiveness of the proposed architectural analysis
system, a comprehensive benchmark suite was designed. The primary
goal of this suite is to objectively measure the analyzer's ability
to correctly extract architectural entities (nodes) and dependency
relationships (edges) across multiple programming languages. This
chapter presents the validation methodology, details the internal
architecture of the benchmark runner, and will subsequently discuss
the aggregated results and provide a qualitative analysis of the
encountered error categories.

\section{Validation Methodology}
\label{sec:validation_methodology}

The validation process is fundamentally based on structural
equivalence: the dependency graph automatically extracted by the
system is programmatically compared against a manually defined
\textit{ground truth} for each benchmark. To ensure robust testing
across paradigms, a dedicated test project was constructed for each
supported language (Rust, Java, Python, C, and C++). These projects
purposefully exercise the core architectural constructs and edge
cases of each language, including module definitions, structured
types (classes, structs, interfaces, traits, enums), function
declarations, inheritance hierarchies, method invocations,
instantiations, and field accesses.

\subsection{The Architectural Oracle (Manifest)}
For each test project, the ground truth is formalized in a YAML
manifest file (typically named \texttt{test.yml}). This manifest acts
as the architectural oracle against which the extracted graph is
verified. The manifest explicitly defines:

\begin{itemize}
  \item \textbf{Target Files}: The list of source code paths that
    constitute the benchmark project, ensuring the analyzer parses
    the exact subset of files required.
  \item \textbf{Expected Nodes}: The architectural entities that the
    system must recognize. Each node specifies a Fully Qualified Name
    (FQN), such as \texttt{models.User.\_\_init\_\_}, and a
    corresponding entity kind (e.g., \texttt{module}, \texttt{class},
    \texttt{struct}, \texttt{function}).
  \item \textbf{Expected Edges}: The semantic dependencies connecting
    the entities. Each edge definition requires a unique identifier,
    a \texttt{source} FQN, a \texttt{sink} (target) FQN, and the
    specific relationship type (e.g., \texttt{calls},
      \texttt{inherits}, \texttt{uses\_type}, \texttt{accesses\_field},
    \texttt{imports}).
\end{itemize}

\subsection{The Benchmark Runner and Validation Logic}
To automate the evaluation process, a dedicated test harness, the
\texttt{benchmark\_runner}, was developed. The runner executes the
full extraction pipeline (Parsing, Structural Extraction, Semantic
Resolution, and Graph Construction) on the test files and
programmatically verifies adherence to the YAML manifest.

Because the analyzer generates a deeply nested, strongly typed
Intermediate Representation (IR), while the YAML manifest uses flat
string identifiers, the validation logic must bridge this structural
gap. The validation algorithm operates through the following steps:

\begin{enumerate}
  \item \textbf{Graph Flattening and Indexing:} To optimize lookups
    and resolve naming mismatches, the runner flattens the extracted
    \texttt{DependencyGraph} into indexing maps. For nodes, it
    traverses the hierarchy of modules and components, extracting the
    nested paths (e.g., \texttt{["models", "User", "\_\_init\_\_"]})
    and joining them into a flat string FQN
    (\texttt{"models.User.\_\_init\_\_"}), associating it with the
    original component instance. Similarly, it constructs an
    adjacency map for the edges, mapping each flattened source string
    to a set of valid sink strings.

  \item \textbf{Node Validation:} The algorithm iterates through the
    \textbf{expected nodes} in the manifest and queries the flattened
    node index. The primary metric is the existence of the node based
    on its FQN. While the system extracts rich typing information
    (e.g., distinguishing between a \texttt{StructuredType(Class)}
    and a \texttt{Function}), the validation relies primarily on
    structural naming matches to maintain compatibility across
    different language paradigms.

  \item \textbf{Edge Validation (Dependency Checks):} For every
    expected edge, the runner performs a two-step validation:
    \begin{itemize}
      \item \textit{Endpoint Verification}: It first ensures that
        both the \texttt{source} and \texttt{sink} nodes were
        successfully extracted. If either endpoint is missing from
        the node index, the edge automatically fails validation due
        to a "Missing Endpoint," isolating the root cause of the error.
      \item \textit{Adjacency Verification}: If both endpoints exist,
        the algorithm queries the adjacency map to determine if a
        dependency relation was established from the source to the sink.
    \end{itemize}
    Currently, the runner focuses on verifying the existence of the
    dependency edge rather than strictly enforcing the specific
    \texttt{DependencyEdgeKind}. However, because the analyzer
    natively extracts highly specific relationship types (such as
    \texttt{UsesParamType} or \texttt{AccessesField}), the
    infrastructure is fully prepared to support strict edge-kind
    validation in future iterations. When an edge fails, the runner
    outputs the list of available sinks for that specific source,
    providing crucial contextual data for debugging the resolution heuristics.
\end{enumerate}

\subsection{Reporting and Metrics}
Upon completing the validation, the runner serializes the results
into both a JSON format (for potential CI/CD automation) and a
human-readable Markdown report. The report dynamically models summary
tables that classify the extraction outcomes into pass/fail statuses.

The evaluation primarily tracks two key metrics:
\begin{itemize}
  \item \textbf{Node Accuracy}: The fraction of expected
    architectural entities successfully identified by the analyzer.
    \[
      \text{Acc}_{nodes} = \frac{|\{n \in N_{expected} \mid n \in
      N_{extracted}\}|}{|N_{expected}|}
    \]
  \item \textbf{Edge Accuracy}: The fraction of expected dependencies
    successfully resolved.
    \[
      \text{Acc}_{edges} = \frac{|\{e \in E_{expected} \mid e \in
      E_{extracted}\}|}{|E_{expected}|}
    \]
\end{itemize}

These metrics provide a quantitative foundation for analyzing the
strengths and limitations of the language-agnostic extraction
heuristics, the results of which are discussed in the following sections.

\section{Standardized Multi-Language Benchmark Results}
\label{sec:benchmark_results}

The standardized benchmark suite comprises test codebases for all five
target languages (C, C++, Java, Rust, and Python), accompanied by
hand-crafted YAML oracles encoding the ground truth for nodes and edges.
The benchmark scenarios specifically assess advanced language mechanisms,
including:
\begin{itemize}
  \item \textbf{C++ Templates}: Class templates, function templates, and
    out-of-line template method definitions.
  \item \textbf{Java Generics}: Parameterized classes, bounded type parameters
    (e.g., \texttt{<T extends Serializable>}), and generic method calls.
  \item \textbf{Rust Generics and Traits}: Generic structs (\texttt{Box<T>}),
    multi-trait bounds (\texttt{T: Display + Clone}), implementation blocks
    for generic types, and \texttt{Deref} coercion.
  \item \textbf{Python Type Variables}: Generic subscriptions (\texttt{List[T]}),
    \texttt{TypeVar} factory calls, and inheritance from \texttt{Generic[T]}.
  \item \textbf{C Structural Constructs}: Header-source file coupling,
    struct field references, and function pointer interactions.
\end{itemize}

Table \ref{tab:benchmark_suite_results} summarizes the accuracy of node
and edge extraction achieved across all five benchmark targets.

\begin{table}[htpb]
  \centering
  \begin{adjustbox}{max width=\textwidth}
  \begin{tabular}{|l|c|c|c|c|c|c|}
    \hline
    \textbf{Language} & \textbf{Exp. Nodes} & \textbf{Ext. Nodes} &
    \textbf{Node Acc.} & \textbf{Exp. Edges} & \textbf{Ext. Edges} &
    \textbf{Edge Acc.} \\ \hline
    \textbf{C}        & 49  & 49  & 100\% & 43 & 43 & 100\% \\ \hline
    \textbf{C++}      & 50  & 50  & 100\% & 44 & 44 & 100\% \\ \hline
    \textbf{Java}     & \multicolumn{3}{c|}{Topology-based oracle} & 46 & 46 & 100\% \\ \hline
    \textbf{Rust}     & 56  & 56  & 100\% & 84 & 84 & 100\% \\ \hline
    \textbf{Python}   & 57  & 57  & 100\% & 58 & 58 & 100\% \\ \hline
  \end{tabular}
  \end{adjustbox}
  \caption{Extraction accuracy across the standardized multi-language benchmark suite.}
  \label{tab:benchmark_suite_results}
\end{table}

In addition to the benchmark suite, an integration test suite comprising 17
dedicated scenarios (\texttt{tests/generics\_test.rs}) validates corner cases,
including empty compilation units, self-referential type assignments, mutual
transitive star-imports, and destructor name preservation. All 17 tests pass
without regressions.

\section{Empirical Evaluation on Large-Scale Open-Source Repositories}
\label{sec:large_scale_evaluation}

To evaluate scalability, throughput, and robustness in real-world environments,
OmniDeps was evaluated against several widely deployed open-source software
repositories exhibiting diverse architectural styles:
\begin{itemize}
  \item \textbf{JUnit 5}: The standard test framework for Java, consisting of
    over 9,600 compilation units. It exercises heavy interface inheritance,
    including complex test extensions implementing up to 21 distinct interfaces
    simultaneously (such as \texttt{KitchenSinkExtension}).
  \item \textbf{Fastjson}: A high-performance Java JSON library with over
    19,000 compilation units, featuring deeply nested classes and pervasive
    generic type parameterizations.
  \item \textbf{ANTLR4}: A multi-language parser generator analyzed across its
    Java and Python runtimes, combining static object hierarchies with dynamic
    package re-exports.
  \item \textbf{JUnit 4}: The predecessor Java framework, analyzed as an
    additional baseline for classic object-oriented hierarchies.
\end{itemize}

Table \ref{tab:large_scale_results} reports the quantitative extraction metrics
and execution times measured on an Apple M-series workstation.

\begin{table}[htpb]
  \centering
  \begin{adjustbox}{max width=\textwidth}
  \begin{tabular}{|l|c|r|r|r|r|r|r|}
    \hline
    \textbf{Repository} & \textbf{Language} & \textbf{Modules} &
    \textbf{Types} & \textbf{Functions} & \textbf{Resolved} &
    \textbf{External} & \textbf{Time} \\ \hline
    \textbf{JUnit 5}    & Java              & 9,622 & 4,102 & 436 & 21,967 & 42,151 & 5.79 s \\ \hline
    \textbf{Fastjson}   & Java              & 19,188 & 6,302 & 0   & 21,179 & 44,517 & 15.20 s \\ \hline
    \textbf{ANTLR4}     & Java / Python     & 4,237 & 1,280 & 613 & 15,834 & 11,299 & 4.80 s \\ \hline
    \textbf{ANTLR4 Runt.} & Java / Python   & 2,053 & 709   & 601 & 2,549  & 10,348 & 2.95 s \\ \hline
    \textbf{JUnit 4}    & Java              & 2,131 & 1,359 & 0   & 3,880  & 7,463  & 1.90 s \\ \hline
  \end{tabular}
  \end{adjustbox}
  \caption{Empirical evaluation on large-scale open-source software repositories.}
  \label{tab:large_scale_results}
\end{table}

\section{Performance and Convergence Analysis}
\label{sec:evaluation_discussion}

The experimental results demonstrate that the language-agnostic architecture
scales effectively to large codebases while maintaining 100\% adherence on
standardized benchmarks. Key operational insights include:

\begin{enumerate}
  \item \textbf{Memoization Impact}: Without the memoization sentinel
    $\mathcal{M}(\Sigma)$, classes implementing extensive interface sets
    (such as \texttt{KitchenSinkExtension} in JUnit 5) caused factorial
    recursion loops ($\mathcal{O}(b!) \approx 2.4 \times 10^{18}$ calls),
    preventing analysis completion. Introducing the sentinel collapsed
    super-scope lookup to $\mathcal{O}(|V_{\text{supers}}| + |E_{\text{supers}}|)$,
    allowing JUnit 5 to complete in 5.79 seconds.
  \item \textbf{Generic Type Preservation}: Evaluating super-types from the
    local scope of the derived class ensures that formal type parameters
    remain bound during base-type resolution (e.g., \texttt{Box<T> extends Base<T>}),
    eliminating false-negative failures.
  \item \textbf{Scope-Aware Cycle Guards}: Scoping query visited keys by
    $\langle \text{ScopeId}, \text{query} \rangle$ successfully prevents
    divergence on self-referential assignments and circular re-exports
    without producing false cycle detections on homonymous identifiers in
    distinct scopes.
\end{enumerate}
